% ---- Introduction -----------------------------------------------------------
% Target length: ~1 page.
% Structure: (1) problem & motivation, (2) gap in prior work,
%            (3) our approach overview, (4) contributions list.

Icons are the visual vocabulary of modern user interfaces. Every button, menu,
notification, and action is accompanied by a small vector graphic that must
communicate its meaning instantly and scale perfectly from 16 to 512 pixels.
Professional icon libraries --- Heroicons, Material Symbols, Tabler, Lucide ---
contain tens of thousands of carefully crafted SVG files, but they cannot cover
the long tail of application-specific concepts: a domain-specific status
indicator, a company-branded element, or a novel UI affordance.

Generating icons programmatically from natural language would allow designers
and developers to describe what they need --- \textit{"a credit card with three
sparkles in the top right corner"} --- and receive a production-ready SVG file
that renders pixel-perfectly at any resolution and remains fully editable in
vector design tools.

\paragraph{The gap in existing approaches.}
Text-to-image diffusion models \citep{rombach2022ldm, saharia2022imagen} can
produce visually plausible icon-like images, but the output is a raster bitmap,
not a vector graphic. Converting a raster to SVG via auto-tracing produces
imprecise, verbose path data unsuitable for production use. Text-to-SVG methods
based on Score Distillation Sampling \citep{jain2023vectorfusion, xing2024svgdreamer}
optimize Bézier control points to match a target image, yielding semantically
reasonable results for illustrations but generating many redundant paths,
lacking the clean single-path-per-shape structure of handcrafted icons, and
taking tens of seconds per image.

LLM-based SVG generation \citep{rodriguez2025starvector, yang2025omnisvg,
xing2025llm4svg} has recently demonstrated that treating SVG as code leads to
more structured, editable outputs. However, these models are trained on general
SVG corpora covering logos, diagrams, and illustrations --- anime-style artwork,
multi-colour paintings, and raster conversions --- with no focus on icons.

More critically, \emph{no existing model targets monochrome icon generation
specifically}. Professional icon sets --- Heroicons, Material Symbols, Tabler,
Lucide --- use a single semantic colour via the \texttt{currentColor} CSS
property, allowing icons to be instantly recoloured to match any brand or design
system. This monochrome constraint is architecturally distinct from multi-colour
SVG generation: the output is structurally simpler (fewer paths, no gradients,
no colour layers) but semantically demanding --- the silhouette alone must convey
the entire meaning of the icon. This focus on \emph{lightweight, pixel-perfect,
immediately brandable monochrome icons} is what sets our work apart from all
prior text-to-SVG studies.

\paragraph{Our approach.}
We fine-tune Qwen3-Coder-30B \citep{qwen3} --- a state-of-the-art
mixture-of-experts code language model --- on a dataset of \num{275912}
(SVG icon, natural language description) pairs using LoRA \citep{hu2022lora}.
The model is selected from a systematic comparison of Qwen and DeepSeek coder
variants (see Figure~\ref{fig:model-comparison}). The natural language
descriptions are generated automatically by a vision language model
(Qwen2.5-VL-7B \citep{qwen25vl}) applied to rendered PNG images of each icon,
following an established automated annotation pipeline
\citep{rodriguez2025starvector, xing2025llm4svg}.

The resulting model, \textbf{\method}, generates valid, renderable, monochrome
SVG icons in under one second. Unlike diffusion-based approaches, the output
is natively editable, consistently structured, uses \texttt{currentColor} for
immediate brand recolouring, and follows the conventions of professional icon
design and development workflows.

\paragraph{Contributions.}
\begin{enumerate}
    \item \textbf{The \dataset{} Corpus.} A curated dataset of \num{275912}
          SVG icons from the Iconify open-source ecosystem, filtered by license
          (excluding ShareAlike and GPL variants), normalized to a consistent
          SVG format, and annotated with VLM-generated natural language
          descriptions. To our knowledge, this is the first publicly documented
          use of Iconify as a machine learning training corpus.

    \item \textbf{An automated icon annotation pipeline.} A reproducible
          pipeline for rendering SVG icons and generating structured
          short and long natural language descriptions using a vision language
          model, with documented prompt engineering and quality metrics.

    \item \textbf{\method.} A monochrome text-to-SVG model fine-tuned with
          LoRA on Qwen3-Coder-30B, trained specifically for single-colour icon
          generation with curriculum ordering (simple to complex) and
          multi-variant instruction augmentation. The monochrome
          \texttt{currentColor} constraint makes outputs immediately usable
          across any brand or design system.

    \item \textbf{An icon-specific evaluation framework.} Metrics tailored to
          icon generation quality: SVG validity rate, rendering fidelity,
          semantic alignment (CLIP score), and structural economy (path count),
          together with a held-out test set of \num{13795} icons for
          reproducible benchmarking.
\end{enumerate}

Code, model weights, and the \dataset{} corpus are released at:
\url{https://github.com/TODO/icon-genai} under open-source licenses.
